% ---------------------------------------------------------------------------
%  QML Lab research project 14: Noise-Aware Training with Error Mitigation
%  Build:  latexmk -pdf main.tex   (needs qmlproject.sty, harvardqml.bst, references.bib)
% ---------------------------------------------------------------------------
\documentclass[11pt,a4paper]{article}
\usepackage{qmlproject}
\qmlheader{Project 14: Noise-aware training}

\hypersetup{pdftitle={QML Lab project 14: Noise-Aware Training with Error Mitigation}}

\begin{document}

\begin{center}
  {\small\textcolor{black!55}{QML Lab \textbullet\ research project 14}}\\[0.5em]
  {\LARGE\bfseries Noise-Aware Training with Error Mitigation}\\[0.6em]
  {\large How much accuracy does zero-noise extrapolation recover for a classifier under realistic noise?}\\[0.6em]
  {\small Living document \textbullet\ last updated \today}
\end{center}

\begin{projectcard}
\small
\begin{tabularx}{\linewidth}{@{}>{\bfseries}l X >{\bfseries}l X@{}}
Quadrant  & \CQ\ (classical data, quantum processor) & Track & Noise and hardware \\
Difficulty & Medium & Tools & Qiskit Aer, Mitiq, PennyLane \\
Builds on & Modules 5 and 8 of the QML Lab & Workspace & \path{qml_projects/noise-aware-training-zne/} \\
\end{tabularx}
\end{projectcard}

\paragraph{Summary.}
On today's hardware, noise shrinks every expectation value a model produces. This project trains and
evaluates the variational classifier of Module 5 under realistic noise models and applies zero-noise
extrapolation (ZNE) to recover noise-free values. It separates two questions that are often mixed up:
whether noise changes the predicted class, and whether it biases the numbers the model relies on for
training and calibration.

% ===========================================================================
\section{Background and key concepts}
% ===========================================================================
\subsection{Noise in variational models}
Gate errors, decoherence and readout errors make an expectation value $E(\lambda)$ depend on the
noise strength $\lambda$, with $\lambda=1$ for the native device. Under global depolarizing noise,
$E(\lambda)\approx(1-p)^{d}E_0$ shrinks exponentially with depth $d$. Noise also flattens the cost
landscape, a noise-induced barren plateau \citep{wang2021noise}. Error mitigation reduces the bias
of expectation values without full error correction \citep{cai2023}.

\subsection{Zero-noise extrapolation}
ZNE \citep{temme2017,kandala2019} measures $E(\lambda_k)$ at amplified noise levels
$1=\lambda_1<\lambda_2<\dots$ and extrapolates to $\lambda=0$. Richardson extrapolation uses
\begin{equation}
  E(0)\approx\sum_k\gamma_kE(\lambda_k),\qquad \sum_k\gamma_k=1,\qquad \sum_k\gamma_k\lambda_k^{j}=0\ \ (j=1,\dots,K-1).
\end{equation}
Digital ZNE amplifies noise by unitary folding, $G\to G(G^\dagger G)^m$, which leaves the ideal circuit
unchanged but multiplies its noise by $\lambda=1+2m$ \citep{giurgica2020}. Mitiq implements these
methods for Qiskit and PennyLane circuits \citep{larose2022}. The price is variance: extrapolated values
need many more shots.

\subsection{Why the class may survive noise}
If noise only rescales $E$ towards 0, the sign of $f(x)=E(x)+b$ changes only for points close to the
decision boundary. Accuracy can therefore be surprisingly robust, while training (which needs gradient
magnitudes) and calibrated probabilities suffer. Measuring both is part of the project.

% ===========================================================================
\section{Research questions}
% ===========================================================================
The project is open-ended: these questions are the starting point, not the finish line. Add new
ones to the README and the journal as they come up.

\begin{rqlist}
  \item Is it better to train with noise, or to train noise-free and mitigate at test time?
  \item How much does ZNE recover in accuracy and in expectation-value error, and at what shot overhead?
  \item Does a device-based noise model lead to different conclusions than simple depolarizing noise?
\end{rqlist}

% ===========================================================================
\section{Method}
% ===========================================================================
\subsection{Data}
Two moons and digits (Project 04 features), with the splits used in Module 5 and Project 04 so that
noise-free results are directly comparable.

\subsection{Models}
The Module 5 re-uploading classifier ($n=2$--4 qubits, $L=1$--3 layers). Noise models in Qiskit Aer:
depolarizing errors on one- and two-qubit gates ($p_1=10^{-3}$, $p_2=10^{-2}$) plus readout error, and a
noise model built from a real device's calibration data. ZNE with scale factors $\{1,3,5\}$ and linear or
Richardson extrapolation through Mitiq.

\subsection{Baselines}
Noise-free training and evaluation (upper bound); noisy evaluation without mitigation; readout-error
mitigation alone.

\subsection{Metrics}
\begin{itemize}
  \item Test accuracy; mean absolute error of $\langle Z\rangle$ against noise-free values.
  \item Signal magnitude $\E|\langle Z\rangle|$ as a function of depth and noise.
  \item Number of circuit executions (shot overhead) per prediction.
\end{itemize}

% ===========================================================================
\section{Planned experiments}
% ===========================================================================
Each experiment becomes one \texttt{[[experiments]]} entry in the workspace, so its runs, parameters
and metrics are stored and compared in the portal.

\begin{table}[h]
\centering\small
\begin{tabularx}{\linewidth}{@{}l>{\raggedright}X>{\raggedright}X>{\raggedright\arraybackslash}X@{}}
\toprule
ID & Experiment & What varies & Main metrics \\
\midrule
  E1 & Noise sweep & $p_2\in[0,0.05]$, no mitigation & accuracy, signal \\
  E2 & ZNE & linear vs Richardson, scale factors & error, overhead \\
  E3 & Training regimes & train clean/noisy, test noisy $\pm$ ZNE & accuracy \\
  E4 & Readout mitigation & on/off & error \\
  E5 & Depth & $L=1..5$ under noise & accuracy, signal \\
\bottomrule
\end{tabularx}
\caption{Experiment plan. Results go into the journal and the results log (\cref{sec:log}).}
\end{table}

% ===========================================================================
\section{Running it in the QML Lab}
% ===========================================================================
The workspace lives in \path{qml_projects/noise-aware-training-zne/}. Scripts import three helpers from
\path{qml_lab/plots.py}:
\begin{itemize}
  \item \texttt{params()} reads the parameters from the portal's run form (defaults when run by hand);
  \item \texttt{save()} stores figures with the run;
  \item \texttt{log\_metrics()} records numbers that appear in the run table.
\end{itemize}

\begin{lstlisting}
qml_projects/noise-aware-training-zne/
  project.toml          metadata + experiments
  README.md             goal, open questions, ideas
  experiments/          runnable scripts
  docs/main.tex         this document
\end{lstlisting}

First experiment: This experiment is \emph{planned}: write the script, add the block below to \texttt{project.toml}, then run \texttt{python manage.py sync\_projects}.

\begin{lstlisting}[language={}]
[[experiments]]
slug = "zne-vqc"
title = "ZNE for the Module 5 classifier under depolarizing noise"
script = "experiments/e01_zne_vqc.py"
figure = "zne_accuracy.png"

[experiments.params]
noise_p1 = 0.001
noise_p2 = 0.01
scale_factors = "1,3,5"
extrapolation = "richardson"
shots = 2000
n_layers = 3
seed = 0
\end{lstlisting}

\begin{lstlisting}[language=bash]
# from the portal folder: pick up changes to project.toml
python manage.py sync_projects
cd qml_projects/noise-aware-training-zne
# run by hand with the default parameters
PYTHONPATH=../../qml_lab python experiments/e01_zne_vqc.py
\end{lstlisting}

% ===========================================================================
\section{Quality bar}
% ===========================================================================
\begin{itemize}
  \item A classical baseline with a similar number of trainable parameters, tuned with the same care \citep{bowles2024}.
  \item At least five random seeds per configuration, reported as mean $\pm$ standard deviation.
  \item How results change with the number of qubits.
  \item At least one run under a realistic noise model.
  \item Report shot overhead next to every mitigated result; mitigation is never free.
\end{itemize}

% ===========================================================================
\section{Open directions}
% ===========================================================================
\begin{itemize}
  \item Probabilistic error cancellation and other mitigation methods \citep{cai2023}.
  \item Noise-aware circuit design: fewer two-qubit gates, hardware-native layouts.
  \item The same study on real hardware (Project 16).
\end{itemize}

% ===========================================================================
\section{Risks and limitations}
% ===========================================================================
\begin{itemize}
  \item ZNE increases variance; small shot budgets can make results worse than unmitigated ones.
  \item Extrapolation assumes a smooth noise response; non-Markovian device noise can break it.
  \item Accuracy alone can hide large biases in expectation values.
\end{itemize}

% ===========================================================================
\section{Results log}
\label{sec:log}
% ===========================================================================
Summarise findings here as they become stable; day-to-day notes belong in the portal journal.

\begin{table}[h]
\centering\small
\begin{tabularx}{\linewidth}{@{}l l >{\raggedright}X >{\raggedright\arraybackslash}X@{}}
\toprule
Date & Run & Finding & Next step \\
\midrule
 & & & \\
 & & & \\
 & & & \\
\bottomrule
\end{tabularx}
\end{table}

\bibliographystyle{harvardqml}
\bibliography{references}

\end{document}
